% Chapter 2, section 5: Acknowledge messages
% Write in plain LaTeX. Code of a guide step: \guide{<step>}; commands: \begin{shell}...\end{shell};
% a screenshot: \screenshot{file}{caption}; a gap in the draft: \placeholder{...}.
% See book/latex/queen.sty for the full list.

\section{Acknowledge messages}

An ack tells the broker what happened to a message. A \texttt{completed} ack moves the consumer
group's cursor past it, and that is all: the message stays in its partition until retention removes
it.

\queenfig{cursor-per-group}{An ack moves the group's cursor; the message stays.}

\why{the cursor is the group's memory of where it is. Keeping the message after the ack is what
lets other groups, and later replays, read it too.}

\subsection{Four answers}

\begin{options}{Status}{Spends a retry}{What happens, and when to use it}
completed & no & Done. The cursor moves past it. \texttt{true} means the same. \\
failed & yes & The lease is released and the message comes back after \texttt{retryDelay}. After the last retry it goes to the dead-letter queue. Use it for failures that may pass. \texttt{false} means the same. \\
retry & no & The lease is released and the message comes back at once. Use it to give work back (a worker that is shutting down). \\
dlq & no & Filed in the dead-letter queue now. Use it when trying again cannot help (a payload that will never parse). \\
\end{options}

A \texttt{failed} or \texttt{dlq} ack can carry the reason (\texttt{error:~'card~declined'} next to
the group), and the dead-letter queue keeps it.

\why{``it did not work'' has different meanings. Sometimes trying again in a while is right,
sometimes now, and sometimes never. Each answer sends the message where it should go.}

\subsection{One message or a whole batch}

\texttt{queen.ack()} takes one message or an array. With an array, one call acks the whole batch.

\guide{pop-ack}

Each message of an array can carry its own answer:

\guide{ack-each-status}

In a batch, the cursor stops at the first message that will come back: a \texttt{completed}
message after it is delivered again too.

\queenfig{batch-ack-rule}{In a batch, the cursor stops at the first message that comes back.}

\why{one call for a hundred messages is much cheaper than a hundred calls. The cursor rule keeps
order: a message is never passed while an earlier one of the same partition is still to be done.}

\subsection{Three things to know}

\begin{itemize}
\item \textbf{Inside \texttt{.each()}, throw to fail a message.} When the handler throws, the client
nacks it and skips the rest of that partition's batch. When the handler acks \texttt{failed} itself
and returns, the client does not know: it goes on with the partition's next messages, their acks are
refused (that part of the batch is coming back anyway), and they are handled twice.
\item \textbf{Every ack answers HTTP 200, even a refused one.} Read \texttt{success} in the answer: a
refused ack usually means the lease ran out and another worker has the message now.
\item \textbf{An ack carries its consumer group and its lease.} The client takes both from the
message. Over HTTP you send them yourself: an ack without its lease is not checked, and can complete
a batch another worker holds.
\end{itemize}

On the site, with the code in all six SDK languages: \href{https://queenmq.com/writing-applications/ack/}{queenmq.com/writing-applications/ack}.
